\subsubsection{Rationale Zahlen}
\label{subsubsec:fundamentals_rational_numbers}

Damit auch Quotienten ganzer Zahlen erfasst werden, erweitert man $\Z$ zu den
\emph{rationalen Zahlen}:
\[
  \Q=\left\{\frac pq\,\middle|\,p\in\Z,\ q\in\N\right\}
  \qquad \text{„Menge aller Brüche aus ganzen Zahlen“}.
\]
\begin{center}
  \begingroup
    \def\NumberSetStage{3}
    \CachedTikzFigure{content/resources/chapters/fundamentals/tikz/number_set_inclusions.tex}
  \endgroup
\end{center}
Jede rationale Zahl lässt sich als Bruch ganzer Zahlen darstellen. Endliche
und periodische Dezimalzahlen sind daher rational, zum Beispiel
\[
  -\frac12=\Cloze{1.8cm}{-0{,}5}\in\Q,
  \qquad
  \frac19=\Cloze{2.0cm}{0{,}\overline{1}}\in\Q.
\]
Die rationalen Zahlen sind bezüglich $+$, $-$, $\cdot$ und der Division durch
eine von null verschiedene Zahl abgeschlossen.

\[
  \begin{aligned}
    \frac12+\left(-\frac14\right)
      &=\Cloze{3.0cm}{\frac14}\in\Q,
    &\qquad
    \frac12-\left(-\frac14\right)
      &=\Cloze{3.0cm}{\frac34}\in\Q,\\[0.5em]
    \frac12\cdot\left(-\frac14\right)
      &=\Cloze{3.0cm}{-\frac18}\in\Q,
    &\qquad
    \frac12:\left(-\frac14\right)
      &=\Cloze{3.0cm}{-2}\in\Q
  \end{aligned}
  \qquad
  \text{mit } \frac12\in\Q,\ -\frac14\in\Q.
\]

Bezüglich des Wurzelziehens sind die rationalen Zahlen allerdings nicht
abgeschlossen. So gibt es keine rationale Zahl, deren Quadrat gleich $2$ ist:
\[
  x^2=2
  \quad\Longrightarrow\quad
  x=\pm\Cloze{1.2cm}{\sqrt2}\notin\Q.
\]

